\begin{ExerciseBlock}
\ExerciseHeading{08}{Lineare Beschleunigung}{linear_kinematics_linearly_increasing_acceleration}
Ein Objekt ruht am Ort $x=x_0$m und erfährt ab $t=0$s die folgende Beschleunigung \[ a(t) = -g + \alpha t \].

\begin{enumerate}
	\item Welche SI-Basiseinheiten haben die Konstanten $g$ und $\alpha$?

	\item Skizzieren Sie den Verlauf der 3 Bewegungsdiagramme ($a$-$t$-Diagramm, $v$-$t$-Diagramm und $s$-$t$-Diagramm) grob.

	\item Bestimmen Sie die Geschwindigkeitsfunktion $v(t)$ und die Ortsfunktion $x(t)$ des Objekts.

\end{enumerate}
\end{ExerciseBlock}
\UseExerciseSolution{linear_kinematics_linearly_increasing_acceleration}
